\section{Sharp cyclicity for all control rates}
\label{sec:curvature-graph}

State derivatives are taken at fixed physical parameters $(p,m,\delta)$,
before substituting the root variables. All reductions remain common to
the fixed family, independently of relative control rates.

An active three-halves term has curvature proportional to \(q^{-1/2}\).
In the both-active outer region, \(q=r^2\) and \(J=rP''\) remove
this singular factor, while \(v=\delta/q\) measures separation relative
to the state. With \(k=R'_{p,0,0}(0)>0\), their first-order balance is
\[
 J=2c_0r+\frac{3k}{4}m-\frac{3kA}{8}v
       +O\bigl((r+|m|+|v|)^2\bigr).
\]
The strength $A\ne0$ makes the curvature equation transverse in $v$.
Exact division then gives a separation graph $\Delta$ with
$\Delta_r=r\mathcal W$; the closing curvature $c_0\ne0$ makes
$\mathcal W$ strictly monotone. These two reductions give a single turn
without choosing a relative parameter rate. If either nondegeneracy is
removed, its corresponding reduction fails.

\subsection{The exact curvature-zero graph}

Let \(g_\nu\) be the actual signed-root companion in
Proposition~\ref{prop:quadratic-passage}. Its dependence on
\((p,m,\delta)\) is suppressed in the notation. Define smooth functions
\begin{equation}\label{eq:root-profiles}
 T_\nu(u)=g_\nu(u)/u,\quad
 \ell_\nu(u)=T_\nu(u)^2+\tfrac{u}{2}(T_\nu^2)'(u),
 \quad H_\nu(u)=\tfrac12\ell_\nu'(u).
\end{equation}
The quotient is filled in smoothly at \(u=0\), with \(T_\nu(0)=1\).
For \(u>0\),
\begin{equation}\label{eq:passage-derivatives}
 \Gamma_\nu'(u^2)=\ell_\nu(u),\qquad
 \Gamma_\nu''(u^2)=u^{-1}H_\nu(u),\qquad
 \ell_\nu(0)=1,\quad H_\nu(0)=\tfrac34 A_\nu.
\end{equation}
These are exact identities obtained by differentiating
\(\Gamma_\nu(u^2)=g_\nu(u)^2\).

For \(r>0\), put
\begin{equation}\label{eq:J-definition}
 q=r^2,\quad \delta=r^2v,\quad
 J(r,m,v,p)=rP_{m,r^2v}''(r^2,p).
\end{equation}
Choose a fixed small \(\eta>0\). All profiles in the following
calculation are evaluated at \(\delta=r^2v\).
Write
\[
 U=T_1(r)^2,\qquad W=U-v,\qquad s=r\sqrt{W}.
\]
After a fixed common shrinking, \(W>1/2\) for \(|v|\le\eta\).
The two excursions are then both active and the exact downstream
penetration is \(y=r^2W>0\). Chain differentiation gives
\begin{align}
 Q&=r^2\{v+WT_2(s)^2\},&
 Q'&=\ell_2(s)\ell_1(r),\label{eq:Q-outer}\\
 Q''&=r^{-1}\{\ell_2(s)H_1(r)
       +W^{-1/2}H_2(s)\ell_1(r)^2\}.\label{eq:Q-curvature}
\end{align}
Consequently
\begin{align}\label{eq:J-exact}
 J={}&rR''(Q)[\ell_2(s)\ell_1(r)]^2\\
 &+R'(Q)\{\ell_2(s)H_1(r)
               +W^{-1/2}H_2(s)\ell_1(r)^2\}.\notag
\end{align}
Every expression on the right is jointly smooth through \(r=0\).
This is smoothness of a root profile; it does not assert ordinary
joint \(C^2\) regularity in \((q,\delta)\).

\begin{lemma}[Root curvature and its zero graph]\label{lem:curvature-profile}
The function \(J\) has a joint smooth extension on a common box
in \((r,m,v,p)\). At its root boundary,
\begin{equation}\label{eq:J-boundary}
 J(0,m,v,p)
 =\tfrac34 R'_{p,m,0}(0)
     \{A_1(p,m,0)+A_2(p,m,0)(1-v)^{-1/2}\}.
\end{equation}
In particular,
\begin{equation}\label{eq:J-transverse}
 J_v(0,0,0,p)=-\tfrac38 R'_{p,0,0}(0)A(p)\ne0.
\end{equation}
There is exactly one smooth graph \(v=V(r,m,p)\), with
\(|V|<\eta\), on which \(J=0\) in this box.
\end{lemma}

\begin{proof}
Formula~\eqref{eq:J-exact} proves smoothness, and
\eqref{eq:passage-derivatives} at \(r=s=0\) gives
\eqref{eq:J-boundary}. Differentiation in \(v\) gives
\eqref{eq:J-transverse}. First choose \(\eta\) so that \(J_v\) has
the strict sign \(-A\) and a uniform nonzero margin on the root
boundary. The signs of \(J(0,0,-\eta,p)\) and
\(J(0,0,\eta,p)\) are opposite. Compactness in the nuisance
parameters permits one further shrinking of \(r,m\) on which these
properties persist throughout the box. The implicit function theorem,
existence by the endpoint signs, and strict monotonicity in \(v\)
give the asserted unique graph.
\end{proof}

The boundary value of the graph is exact:
\begin{equation}\label{eq:V-boundary}
 V(0,m,p)=1-\left(\frac{-A_2(p,m,0)}{A_1(p,m,0)}\right)^2,
 \qquad V_m(0,0,p)=\frac{2}{A(p)}.
\end{equation}
Indeed the ratio equals \(1-m/A_1(p,m,0)\); both individual strengths
retain their dependence on \(m\).

To limit this graph to one turn on a fixed interval, divide first along
the cancelling slice and then in the independent strength variable.
Finite cancellation gives
\begin{equation}\label{eq:Q-aligned}
 Q_{0,0}(q,p)=q+q_+^{5/2}L(\sqrt{q_+},p).
\end{equation}
For positive \(q=r^2\) it follows, by ordinary state differentiation,
that
\begin{equation}\label{eq:center-curvature}
 P_{0,0}''(r^2,p)=2c_0(p)+rN(r,p),\qquad
 c_0(p)=\tfrac12R_{p,0,0}''(0),
\end{equation}
where \(N\) is jointly smooth. The offset and linear jets are allowed
to vary in \(p\) throughout this identity.

Let \(V_0(r,p)=V(r,0,p)\). Since \(J(r,0,V_0,p)=0\), the mean-value
identity in the independent variable \(v\) yields
\[
 0=r\{2c_0+rN\}
       +V_0\int_0^1J_v(r,0,tV_0,p)\,dt.
\]
The integral has a nonzero common margin. Therefore \(V_0=rB(r,p)\)
with \(B\) smooth and
\begin{equation}\label{eq:B-boundary}
 B(0,p)=\frac{16c_0(p)}{3R'_{p,0,0}(0)A(p)}.
\end{equation}
Hadamard division in \(m\) now gives the exact identity
\begin{equation}\label{eq:V-division}
 V(r,m,p)=mH(r,m,p)+rB(r,p),\qquad
 H(0,0,p)=\frac{2}{A(p)}.
\end{equation}
The profile $H$ retains the $m$-dependence of $R$ and both passages.

Define the physical separation graph
\begin{equation}\label{eq:Delta-definition}
 \Delta(r,m,p)=r^2V(r,m,p).
\end{equation}
Its derivative factors exactly:
\begin{equation}\label{eq:Delta-derivative}
 \Delta_r=r\mathcal W,\qquad \mathcal W=2V+rV_r.
\end{equation}
The boundary coefficients in \eqref{eq:V-division} give
\begin{equation}\label{eq:W-boundary}
 \mathcal W_m(0,0,p)=\frac4{A(p)},\qquad
 \mathcal W_r(0,0,p)=3B(0,p)
   =\frac{16c_0(p)}{R'_{p,0,0}(0)A(p)}\ne0.
\end{equation}
Joint smoothness and compactness in \(p\) give a common smaller box
on which \(\mathcal W_r\) has a uniform nonzero margin with sign
\(c_0/A\).

\begin{proposition}[One nondegenerate turn]\label{prop:one-fold-graph}
For fixed \((m,p)\), the graph \(r\mapsto\Delta(r,m,p)\) has at most
one positive critical point on a common interval. Any such point
is nondegenerate. In the outer region \(|\delta|/r^2\le\eta\),
\begin{equation}\label{eq:curvature-factor}
 P_{m,\delta}''(r^2,p)
 =r^{-3}\mathcal U(r,m,\delta/r^2,p)
                  [\delta-\Delta(r,m,p)],
\end{equation}
where \(\mathcal U\) is jointly smooth and nonzero, of sign \(-A\).
For fixed controls the ordinary outer curvature zeros have total
multiplicity at most two.
\end{proposition}

\begin{proof}
For fixed \((m,p)\), \(\mathcal W\) is strictly monotone and has
at most one positive zero. At that zero,
\(\Delta_{rr}=r\mathcal W_r\ne0\). Hence \(\Delta\) has at most
two monotone branches.

For the factorization, define
\[
 \mathcal U(r,m,v,p)
  =\int_0^1J_v(r,m,V+t(v-V),p)\,dt .
\]
The derivative bound from Lemma~\ref{lem:curvature-profile}
gives its sign and nonvanishing. The exact mean-value identity is
\(J=(v-V)\mathcal U\). Substituting
\(v=\delta/r^2\) in \eqref{eq:J-definition} gives
\eqref{eq:curvature-factor}.
A horizontal line meets a one-turn graph at at most two transverse
points, or at its sole turn. Transverse intersections are simple;
an intersection at the nondegenerate turn has order two.
Changing from \(r>0\) to \(q=r^2\) is a smooth local diffeomorphism,
and the prefactor in \eqref{eq:curvature-factor} never vanishes.
This proves the multiplicity assertion as well as the distinct count.
\end{proof}

We now cover the complement of the outer region, where an active input
may approach its event at an unrelated rate.

\subsection{The common-interval count}\label{sec:inner}

Use the actual thresholds $0$ and $a_\delta=\Gamma_1^{-1}(\delta)$
from \eqref{eq:passage:true-thresholds}, with the active intervals determined
there. All constants are uniform on the fixed parameter reduction.

In the both-active event layer, compare the difference of the singular
curvatures with their sum. Separation keeps the active inputs at a fixed
relative distance; the difference then dominates both the relative error
and the bounded closing curvature.

Put \(h=|\delta|\), \(\sigma=\sgn\delta\), and let \(L\ge2\) be fixed.
The first-strength observations give, uniformly,
\[
 A_1(p,m,\delta)=A(p)+O(|m|+h),\qquad
 A_2(p,m,\delta)=-A(p)+O(|m|+h).
\]
The smooth root profiles in \eqref{eq:passage:quadratic-jet} yield
\begin{align}
 \Gamma_\nu'(x)
  &=1+\tfrac32 A_\nu\sqrt{x}+O(x),\label{eq:inner-derivative}\\
 \Gamma_\nu''(x)
  &=\tfrac34 A_\nu x^{-1/2}+O(1),\qquad x>0.\label{eq:inner-curvature}
\end{align}
Differentiating the smooth root profiles gives these remainders and their
required derivative bounds. Each $\Gamma_\nu'$ is uniformly positive and
bounded, with active curvature of the sign of its strength.

For \(0<q\le Lh\),
\begin{equation}\label{eq:first-correction}
 |\Gamma_1(q)-q|\le Cq^{3/2}\le C L^{3/2}h^{3/2}.
\end{equation}
After reducing \(h\), this is at most \(h/2\).
The inverse derivative estimates give \(h/2<a_h<2h\).
Every positive penetration in this layer is at most \((L+2)h\).

\begin{lemma}[Curvature signs in the entire event layer]\label{lem:event-layer}
For every fixed \(L\ge2\), one may choose common sufficiently small
bounds on \(|m|\) and \(h=|\delta|\) such that, for \(0<|\delta|\),
the actual curvature signs are
\begin{equation}\label{eq:event-signs}
 \sgn P''=
 \begin{cases}
 \sgn c,&q<\min(0,a_\delta),\\
 \sgn(\delta A),&\min(0,a_\delta)<q<\max(0,a_\delta),\\
 \sgn(-\delta A),&
       q>\max(0,a_\delta),\quad q\le L|\delta|.
 \end{cases}
\end{equation}
The last interval is both-active. There is no curvature zero in any
of these open regions. In the active regions
\(|Q''|\ge\kappa_L h^{-1/2}\) for a uniform \(\kappa_L>0\).
All estimates apply arbitrarily close to the actual events.
\end{lemma}

\begin{proof}
Below both thresholds both passages are inactive, so \(Q=q\) and
\(P''=R''\), with sign \(c\).

If \(\delta=h>0\), the one-active interval is \(0<q<a_h\), where
\(Q=\Gamma_1(q)\). Its curvature has sign \(A\), and
\eqref{eq:inner-curvature} gives
\(|Q''|\ge\kappa(2h)^{-1/2}\) after shrinking.
If \(\delta=-h<0\), that interval is \(-h<q<0\), where
\(Q=-h+\Gamma_2(q+h)\). Its curvature has sign \(-A\) and the same
uniform lower bound. These formulas and estimates hold uniformly as either endpoint is approached from within the open one-active interval.

On the both-active part put
\[
 y=\Gamma_1(q)-\delta>0.
\]
The exact chain rule is
\[
 Q''=\Gamma_2'(y)\Gamma_1''(q)
          +\Gamma_2''(y)\Gamma_1'(q)^2.
\]
Equations~\eqref{eq:inner-derivative}--\eqref{eq:first-correction}
therefore give
\begin{equation}\label{eq:relative-inner-error}
 Q''=\tfrac34 A(q^{-1/2}-y^{-1/2})+E,\qquad
 |E|\le C_L(\sqrt h+h+|m|)
                 (q^{-1/2}+y^{-1/2}).
\end{equation}
Each error is measured relative to the sum of the two actual singular
weights, so the formula remains uniform as either positive argument tends
to zero.

By \eqref{eq:first-correction},
\(\sigma(q-y)\ge h/2\) and \(q,y\le(L+2)h\). Hence
\[
 -\sigma\frac{q^{-1/2}-y^{-1/2}}{q^{-1/2}+y^{-1/2}}
 =\frac{|q-y|}{(\sqrt q+\sqrt y)^2}
 \ge\frac1{8(L+2)}=:d_L>0.
\]
First reduce \(|m|\), and then \(h\), so the error in
\eqref{eq:relative-inner-error} is less than half the principal bound.
Thus \(\sgn Q''=\sgn(-\sigma A)\) and
\(|Q''|\ge\kappa_Lh^{-1/2}\), uniformly up to either event.

Finally the exact closing-leg chain rule is
\[
 P''=R''(Q)(Q')^2+R'(Q)Q''.
\]
The first term is uniformly bounded, whereas \(R'\) has a positive
margin and the preceding active curvature bound diverges as
\(h\downarrow0\). A further fixed reduction proves
\eqref{eq:event-signs}.
\end{proof}

\begin{proposition}[The full distinct count]\label{prop:common-count}
On one common sufficiently small physical interval, every return in
Theorem~\ref{thm:six-main} has at most six distinct fixed points.
\end{proposition}

\begin{proof}
First choose the outer box and its fixed \(\eta\) in
Lemma~\ref{lem:curvature-profile}. Set \(L=1/\eta\ge2\), shrinking
\(\eta\) if needed. Apply Lemma~\ref{lem:event-layer} with this one
\(L\): reduce the common strength bound and then the separation bound.
All further reductions are made inside the assumed common actual
first-return domain. Choose \(q_0,\delta_0\) so both actual thresholds remain
strictly inside \(I=(-q_0,q_0)\), for example using
\(\delta_0<q_0/4\) and the uniform inverse derivative bound.

For \(\delta\ne0\), there are no ordinary curvature zeros on the
three event regions in \eqref{eq:event-signs}. Any remaining curvature
zero must have \(q>L|\delta|>0\), so it is both-active with
\(|\delta|/q<\eta\). Proposition~\ref{prop:one-fold-graph} exhausts
this region and gives at most two ordinary curvature zeros.
The singularity at each distinct threshold has a nonzero active
coefficient, so it is not an ordinary curvature zero.

The derivative \(P'\) is continuous at the thresholds.
It is strictly monotone on at most five consecutive pieces: two
threshold junctions and at most two ordinary curvature zeros
are the only possible divisions. If the signs of the adjacent
curvatures agree, monotonicity extends through the threshold;
if an outer curvature zero does not change sign, it adds no new
piece. There is no interval of zero curvature. Thus \(P'-1\)
has at most five distinct zeros. Rolle's theorem applies to the
ordinary \(C^1\) displacement \(D=P-q\), including its event values,
and gives at most six distinct zeros.

The aligned case follows from the next corollary. The interval and control
bounds are independent of all ratios, since the event-layer estimate holds
uniformly up to either event.
\end{proof}

\begin{corollary}[The four-point slices]\label{cor:four-slices}
On the same interval the distinct fixed-point bound is four when
\(\delta=0\) or \(m=0\), and two when \(m=\delta=0\).
\end{corollary}

\begin{proof}
For \(\delta=0\), a positive curvature zero has \(V(r,m,p)=0\).
Equations~\eqref{eq:B-boundary}--\eqref{eq:V-division} give
\[
 V_r(0,0,p)=B(0,p)=\frac{16c_0(p)}{3k(p)A(p)}\ne0.
\]
Compactness in \(p\) gives a common nonzero margin for \(V_r\)
on a smaller box. Thus, for fixed \(m,p\), there is at most one
positive zero of \(V\), and it is simple. Formula~\eqref{eq:curvature-factor}
and the local diffeomorphism \(q=r^2\), \(r>0\), preserve its ordinary
order. Together with at most one sign change at the aligned threshold,
this gives at most three monotonicity pieces for \(P'\) and four fixed points.

For \(m=0,\delta\ne0\), the outer separation graph is
\(\Delta=r^3B(r,p)\), with sign \(c_0/A\) and nonzero derivative.
It has at most one curvature intersection, necessarily simple, and
one is possible only if \(c_0\delta A>0\).
When this inequality holds, the first two signs in
\eqref{eq:event-signs} agree, leaving one threshold change and
at most one outer change. If \(c_0\delta A<0\), there is no outer
intersection and at most two threshold changes. Again \(P'\)
has at most three monotonicity pieces.

At \(m=\delta=0\), equation~\eqref{eq:center-curvature} and the
inactive closing map give the same strict sign \(c_0\) throughout
the interval. The displacement is strictly convex or strictly
concave, including across \(q=0\), and has at most two zeros.
\end{proof}

Proposition~\ref{prop:field-recognition} and
Lemma~\ref{lem:auxiliary-controls} transfer all four bounds to dependent
field observations by restriction. Propositions~\ref{prop:single-scale-six}
and~\ref{prop:sharp-slice-attainment} below attain them in every full-rank
actual family.

In the transported flow boxes, the contact inputs to this criterion are
the two values \(\beta_\nu(0,0)\), by
\eqref{eq:passage:strength-jet}. In an initial second-port coordinate,
the second strength is multiplied by \(\sqrt{S_1'(\delta)}\)
under the transport of Proposition~\ref{prop:field-recognition}.
Thus cancellation is a condition on the local branch jets and the
regular connector derivative, in addition to the two regular-return
jets and its nonzero curvature.


\subsection{The cubic critical value}

Write \(k(p)=R'_{p,0,0}(0)>0\), with \(p\) in the fixed compact
parameter reduction. Its coordinates may include the offset and
multiplier controls.

\begin{proposition}[The cubic outer-curvature window]
\label{prop:cubic-curvature-window}
After a common further reduction, the signed critical-point graph is
\[
 \rho(m,p)=-\frac{k(p)m}{4c_0(p)}+O(m^2),\qquad \rho(0,p)=0.
\]
It is smooth through zero. The positive critical point exists exactly
when \(mc_0(p)<0\), and then equals \(r_{\rm c}=\rho(m,p)\).
Its critical value has an exact smooth division
\begin{equation}\label{eq:cubic-critical-value}
 \delta_{\rm c}(m,p)=\Delta(\rho(m,p),m,p)=m^3\kappa(m,p),
 \qquad \kappa(0,p)=\frac{k(p)^2}{24A(p)c_0(p)^2}.
\end{equation}
The remainder and all reductions are uniform in \(p\).
On a fixed interval \(0<r<\sqrt{q_0}\), reduced as below, the
physical return has exactly two simple outer curvature zeros if and
only if
\begin{equation}\label{eq:curvature-window}
 mc_0(p)<0,\qquad 0<\delta/\delta_{\rm c}(m,p)<1.
\end{equation}
At \(\delta=\delta_{\rm c}\), \(mc_0<0\), it instead has one
ordinary double zero of the outer curvature. This curve imposes neither
\(D=0\) nor \(D'=0\), and is not a fold curve of periodic orbits.
\end{proposition}

\begin{proof}
Apply the implicit function theorem to \(\mathcal W=0\).
The boundary derivatives \eqref{eq:W-boundary} and the uniform margin
give a unique smooth signed solution \(\rho=m\gamma\), with
\(\gamma(0,p)=-k/(4c_0)\). Its positive part exists precisely for
\(mc_0<0\).
Substitution into \(\Delta=r^2(mH+rB)\) gives the smooth quotient
\[
 \kappa=\gamma^2\{H(m\gamma,m,p)+\gamma B(m\gamma,p)\}.
\]
The boundary coefficients yield \(\gamma B=-4/(3A)\), hence
\eqref{eq:cubic-critical-value}. Its sign is \(1/A\); in the
positive-critical-point sector the critical value has sign
\(-c_0/A\).

First fix \(R=\sqrt{q_0}\) inside the common root box. Then reduce
\(|m|\), uniformly in \(p\), so \(0<r_{\rm c}<R/2\) when
\(mc_0<0\), and
\[
 |mH(R,m,p)|<\tfrac12R|B(R,p)|.
\]
Thus \(\Delta(R,m,p)\) has sign \(c_0/A\), opposite to the
nonzero critical value. Also require \(|\delta_{\rm c}|\) to lie
within the common separation bound. The unique turn and
\(\Delta(0)=0\) give exactly two transverse intersections with the
line \(\delta\) precisely for \eqref{eq:curvature-window}.
At each intersection \(|\delta|/r^2=|V(r,m,p)|<\eta\), so both
belong to the genuine outer chart. Conversely two distinct outer
intersections force the unique positive turn between them and place
\(\delta\) strictly between zero and its critical value.
The factorization \eqref{eq:curvature-factor} transfers their ordinary
orders to \(P''\). At the critical value, the nondegenerate turn gives
\(D''=D'''=0\), \(D''''\ne0\), without conditions on \(D,D'\).
\end{proof}

Figure~\ref{fig:curvature-graph} illustrates a permitted one-turn shape.
\input{figures/curvature-graph}


\subsection{Sharpness in full-rank families}\label{subsec:single-scale-sharpness}

Fix a smooth physical observation inverse
\(\lambda=\Lambda(a,b,m,\delta;\xi_0)\) provided by \eqref{eq:rank},
with \(\xi_0\) fixed, and put
\(\lambda_0=\Lambda(0,0,0,0;\xi_0)\). All passage profiles below
are evaluated at this same actual inverse. Write
\[
 A_0=A_1(\lambda_0),\quad
 c_0=\tfrac12R_{\lambda_0}''(0),\quad
 B=|A_0|,\quad C=|c_0|,
\]
\[
 s=-\sgn c_0,\qquad \tau=\sgn(sA_0).
\]
We first allow an arbitrary scaled strength coefficient, including zero:
\begin{equation}\label{eq:weighted-general-controls}
 \lambda_\eps=\Lambda(\alpha\eps^4,\beta\eps^2,
                         M\eps,d\eps^3;\xi_0),\qquad
 \theta=(\alpha,\beta,M,d),\qquad \eps>0.
\end{equation}

\begin{lemma}[Weighted actual-return division]\label{lem:weighted-quartic}
Let \(K\Subset(0,\infty)\), \(K_-\Subset(-\infty,0)\), and let
\(\theta\) range over a fixed compact set. For sufficiently small
\(\eps>0\), the actual normalized displacements satisfy
\begin{align}
 \mathcal F_\eps(z)&:=\eps^{-4}D_{\lambda_\eps}(\eps^2z^2)
 =F_\theta(z)+\eps E_+(\eps,z,\theta),\label{eq:quartic-limit}\\
 F_\theta(z)&=\alpha+\beta z^2+Mz^3+c_0z^4+\tfrac32A_0dz,\notag\\
 \mathcal G_\eps(x)&:=\eps^{-4}D_{\lambda_\eps}(\eps^2x)
 =G_\theta(x)+\eps E_-(\eps,x,\theta),
 \label{eq:quartic-negative-limit}\\
 G_\theta(x)&=\alpha+\beta x+c_0x^2.\notag
\end{align}
The remainders have auxiliary jointly smooth extensions to signed
\(\eps\) near zero.
For every prescribed finite integer \(k\), uniformly on these compacts,
\begin{equation}\label{eq:quartic-mixed-norm}
 \max_{j+|\gamma|\le k}
 \left\|\partial_z^j\partial_\theta^\gamma
                 (\mathcal F_\eps-F_\theta)\right\|_\infty\le C_k\eps,
\end{equation}
and similarly with \(x,\mathcal G_\eps\). These estimates include
compact control sets meeting \(M=0\). Derivatives in \(\eps\) are
bounded on the smooth extensions; they are not included in the \(O(\eps)\) norm in
\eqref{eq:quartic-mixed-norm}.
\end{lemma}

\begin{proof}
Use the exact identities of Lemma~\ref{lem:exact-value-division}.
Substitute \(r=\eps z\), \(v=d\eps/z^2\) into
\eqref{eq:quartic-Hadamard}, retaining the profile dependence on
\(\delta=r^2v=d\eps^3\) and all other actual controls. Exactly,
\[
 Q-q=\eps^4V,\qquad
 V=Mz^3+dzH+\eps z^5L_0+\tfrac34M^2\eps^2z^4.
\]
With \(Y=z^2+\eps^2V\), the exact closing expansion gives
\begin{equation}\label{eq:quartic-exact-residual}
 \mathcal F_\eps=
 \alpha+\beta Y+V+c(\lambda_\eps)Y^2
                   +\eps^2Y^3S(\eps^2Y,\lambda_\eps).
\end{equation}
Since \(z\) stays away from zero, all arguments remain in the common
root box, whose positive square-root factor stays bounded away from zero.
At \(\eps=0\), \(H=3A_0/2\) and \(c=c_0\), giving \(F_\theta\).
Taylor integral division in \(\eps\) proves the joint remainder and
every finite mixed estimate in \eqref{eq:quartic-mixed-norm}, without
dividing by \(M\). State differentiation holds \((\eps,\theta)\), hence
\(\lambda_\eps\), fixed: the varying value \(v=d\eps/z^2\) keeps
the physical separation \(r^2v=d\eps^3\) fixed.
The signed extension uses \(r=\eps z\), not the physical
root \(|\eps|z\); physical equality is asserted for \(\eps>0\).

For negative \(x\) in a fixed compact set, \(q=\eps^2x\) is below
both events for small positive scale, since \(\delta=O(\eps^3)\) and
\(|a_\delta|\le2|\delta|\). Thus \(Q=q\) exactly, and
\eqref{eq:quartic-closing-division} gives
\[
 \mathcal G_\eps=\alpha+\beta x+c(\lambda_\eps)x^2
                   +\eps^2x^3S(\eps^2x,\lambda_\eps).
\]
This proves the negative formula and its joint smooth extension.
\end{proof}

For the six-point construction specialize to
\(\alpha=s\alpha_4\), \(\beta=s\bsc\), \(M=s\) and
\(d=\tau\dsc\), giving
\begin{equation}\label{eq:quartic-weights}
 m=s\eps,\qquad \delta=\tau \dsc\eps^3,\qquad
 b=s\bsc\eps^2,\qquad a=s\alpha_4\eps^4,\qquad \eps>0.
\end{equation}
Here \(\alpha_4,\bsc,\dsc\) are scaled controls. The hats distinguish them from the
ecological parameters $\beta,d$ used in Section~\ref{sec:harvesting}.
Multiplying the normalized displacements by \(s\), and using
\(s\tau A_0=B\), \(sc_0=-C\), gives
\(\mathcal F_{4,\eps}:=s\mathcal F_\eps=F_0+O(\eps)\) and
\(\mathcal G_{4,\eps}:=s\mathcal G_\eps=G_0+O(\eps)\), where
\[
 \begin{aligned}
 F_0(z)&=\alpha_4+\bsc z^2+z^3-Cz^4+\tfrac32B\dsc z,\\
 G_0(x)&=\alpha_4+\bsc x-Cx^2.
 \end{aligned}
\]
These expansions have all the mixed estimates of
Lemma~\ref{lem:weighted-quartic}.

Take distinct positive numbers \(r_i\), \(1\le i\le4\), with
\(\sum r_i=1/C\), and let \(e_j\) be their elementary symmetric
functions. Choosing
\begin{equation}\label{eq:quartic-Vieta}
 \alpha_4=-Ce_4,\qquad \bsc=-Ce_2,\qquad \dsc=\frac{2C}{3B}e_3
\end{equation}
gives simultaneously
\[
 F_0(z)=-C\prod_{i=1}^4(z-r_i),\qquad
 G_0(x)=-C(x^2+e_2x+e_4).
\]
The six pair products have product \(e_4^3\); AM--GM yields
\(e_2\ge6\sqrt{e_4}\). Thus \(e_2^2-4e_4>0\), and the quadratic
has two distinct negative roots. This relation pairs the positive
quartet with two inactive companions.

\begin{proposition}[Single-scale physical sharpness]
\label{prop:single-scale-six}
Every actual family satisfying \eqref{eq:rank} attains six simple fixed
points on the state scale $\eps^2$, uniformly on an open scaled-control
set. Their stability alternates: three are attracting and three repelling.
\end{proposition}

\begin{proof}
Choose any distinct positive $r_i$ with sum $1/C$ and the controls
\eqref{eq:quartic-Vieta}. The four roots of $F_0$ and the two negative
roots of $G_0$ have strict derivative margins. Lemma~\ref{lem:weighted-quartic}
continues them on six disjoint isolators, uniformly over a small open
scaled-control set. The same joint estimates allow \(M\) to vary near
\(s\), so this set may be taken open in all four scaled controls
\((\alpha,\beta,M,d)\). The observation inverse realizes these controls in the
actual family; fixed scaled controls give a smooth physical path tending
to the reference. The common upper bound excludes additional fixed points.
At the displayed controls, for negative states $D_q=s\eps^2G_0'(x)+O(\eps^3)$, and for positive
states $D_q=s\eps^2F_0'(z)/(2z)+O(\eps^3)$. In state order these signs
are $s(+,-,+,-,+,-)$. Since the return is increasing and close to the
identity, they give three attracting and three repelling hyperbolic cycles;
these signs persist on the open scaled-control set.
\end{proof}

\begin{proposition}[Sharpness on the exceptional slices]
\label{prop:sharp-slice-attainment}
Every actual family satisfying \eqref{eq:rank} attains four simple
fixed points on \(\delta=0\), four on \(m=0\), and two on
\(m=\delta=0\). Each attainment occurs in the common interval \(I\)
for every sufficiently small \(\eps>0\) along a smooth physical
parameter path tending to the reference, and persists on an open
scaled-control set relative to the corresponding slice.
\end{proposition}

\begin{proof}
Use the weights \eqref{eq:weighted-general-controls}. The following table
specifies the paths and the limiting simple roots; positive roots use the
coordinate $z$, and negative roots the coordinate $x$:
\begin{equation}\label{eq:sharp-slice-controls}
 \begin{array}{c|c|c|c}
 \text{slice}&(\alpha,\beta,M,d)/(sC)&z>0&x<0\\ \hline
 \delta=0&(36/11,-85/11,60/11,0)&1,2,3&x_\delta\\
 m=0&(36,25,0,-40/A_0)&1,2,3&x_m\\
 m=\delta=0&(1,0,0,0)&1&-1
 \end{array}
\end{equation}
Here $x_\delta=(-85-\sqrt{8809})/22$ and
$x_m=(25-\sqrt{769})/2$. For the first two rows, direct substitution in
Lemma~\ref{lem:weighted-quartic} gives respectively
\[
 \begin{aligned}
 F_\theta/(sC)&=-(z-1)(z-2)(z-3)(z+6/11),&
 G_\theta/(sC)&=36/11-85x/11-x^2,\\
 F_\theta/(sC)&=-(z-1)(z-2)(z-3)(z+6),&
 G_\theta/(sC)&=36+25x-x^2.
 \end{aligned}
\]
Only positive $z$ roots represent physical states in this chart; the
negative companions come from $G_\theta$. The mixed estimates include
$M=0$.

For the third row, \eqref{eq:center-displacement} gives, on compact
$x$ intervals including zero,
\[
 \frac{D_{\lambda_\eps}(\eps^2x)}{sC\eps^4}
 =1-x^2+O_{C^1}(\eps).
\]
Indeed, with $p_\eps=(sC\eps^4,0,\xi_0)$, its root remainder is
$\eps x_+^{5/2}\widetilde L(\eps\sqrt{x_+},p_\eps)/(sC)$.
This term and its state derivative are uniformly $O(\eps)$, including
at zero; the cubic remainder is $O(\eps^2)$, and
$c_0(p_\eps)=c_0+O(\eps^4)$.

Choose disjoint compact isolators about the listed roots, with strict
endpoint-sign and derivative margins. These margins persist on relative
open scaled-control sets with compact closures: $d=0$ in the first row,
$M=0$ in the second, and $M=d=0$ with $(\alpha,\beta)$ near $(sC,0)$
in the third. The actual thresholds satisfy
$|a_\delta|/\eps^2\le2|d|\eps$ and $0/\eps^2=0$, so every positive
isolator is both-active and every negative isolator lies below both events
for small scales, uniformly on those sets. One scale reduction puts all
controls in the fixed observation chart and all scaled isolators in $I$.
Joint $C^1$ convergence continues one simple root in each isolator;
Corollary~\ref{cor:four-slices} excludes every other root in $I$, including
at the events. The fixed observation inverse gives smooth physical paths
and the asserted relative persistence.
\end{proof}
